% !TEX root = task-1-15-2.tex
% Editable task bank. Append new BankTask blocks at the end of the document.
\documentclass[12pt,a4paper]{article}
\usepackage{physics-textbook}
\usepackage{graphicx,needspace}
\newif\ifShowSolutions
\ifdefined\TutorEdition
  \ShowSolutionstrue
\else
  \ShowSolutionsfalse % Change to \ShowSolutionstrue to display solutions.
\fi
\newcounter{banktask}
\newcommand{\BankTask}[3]{%
  \refstepcounter{banktask}%
  \Needspace{35mm}\par\medskip\noindent\begin{minipage}{\linewidth}%
  \noindent\textbf{\thebanktask.}\quad #2\par
  \ifShowSolutions
    #3\par
    \ifdefined\TutorEdition
      \ifcsname TutorSource#1\endcsname
        \smallskip{\small\csname TutorSource#1\endcsname\par}%
      \fi
    \fi
  \fi
  \end{minipage}\par\medskip
}
\begin{document}
\BookSubsection{1.15.2}{Мощность}


% Source sdamgia-545: https://phys-ege.sdamgia.ru/problem?id=545
\BankTask{sdamgia-545}{
Под действием силы тяги в $1000\,\text{Н}$ автомобиль движется с постоянной скоростью $72\,\text{км/ч}$. Какова мощность двигателя?
}{
\smallskip\noindent\textbf{Решение}\par
Переведем скорость в систему СИ: $72\,\text{км/ч}=20\,\text{м/с}$ Мощность двигателя связана со скоростью движения и силой тяги соотношением: 


$N=F\cdot v=1000\,\text{Н}\cdot20\,\text{м/с}=20\,\text{кВт}.$

Ответ: $20$.
}

% Source sdamgia-525: https://phys-ege.sdamgia.ru/problem?id=525
\BankTask{sdamgia-525}{
Лебедка равномерно поднимает груз массой $200\,\text{кг}$ на высоту $3\,\text{м}$ за $5\,\text{с}$. Какова мощность двигателя лебедки? Ускорение свободного падения принять равным $10\,\text{м/с}^2$.
}{
\smallskip\noindent\textbf{Решение}\par
Поскольку лебедка поднимает груз равномерно, по второму закону Ньютона сила, с которой тянет лебедка, в точности равна по модулю силе тяжести, действующей на груз. Таким образом, лебедка совершает работу: 


$A=Fh=mgh=200\,\text{кг}\cdot10\,\text{м/с}^2\cdot3=6000\,\text{Дж}.$
 
Следовательно, мощность двигателя лебедки равна: 


$N=\frac At=\frac{6000\,\text{Дж}}{5\,\text{с}}=1200\,\text{Вт}.$

Ответ: $1200$.
}

% Source sdamgia-6812: https://phys-ege.sdamgia.ru/problem?id=6812
\BankTask{sdamgia-6812}{
Телу массой $0{,}2\,\text{кг}$ сообщили вертикально направленную начальную скорость $10\,\text{м/с}$. Пренебрегая сопротивлением воздуха, определите модуль средней мощности силы тяжести, действовавшей на тело во время подъема до максимальной высоты. Ускорение свободного падения считать равным $10\,\text{м/с}^2$.
}{
\smallskip\noindent\textbf{Решение}\par
Мощность силы равна скалярному произведению силы на скорость. Поскольку на тело действует постоянная сила, сила тяжести, а скорость убывает линейно, то среднюю мощность можно определить как полусумму начальной и конечной мощностей. Примем во внимание, что в точке наивысшего подъема скорость тела равна нулю:


$|N_{\text{ср}}|=\left|\frac{N_{\text{нач}}+N_{\text{кон}}}{2}\right|=\frac{mgv_{\text{нач}}+mgv_{\text{кон}}}{2}=$
$=\frac{v_{\text{нач}}+v_{\text{кон}}}{2}\cdot mg=\frac{10\,\text{м/с}+0\,\text{м/с}}2\cdot0{,}2\,\text{кг}\cdot10\,\text{м/с}^2=10\,\text{Вт}.$



Ответ: $10$.
}

% Source sdamgia-10308: https://phys-ege.sdamgia.ru/problem?id=10308
\BankTask{sdamgia-10308}{
Тело массой $1\,\text{кг}$ свободно падает на землю с высоты $20\,\text{м}$. Какую среднюю мощность развивает сила тяжести за время падения тела?
}{
\smallskip\noindent\textbf{Решение}\par
Тело падает в поле тяжести Земли, высота зависит от времени следующим образом:


$h=\frac{gt^2}{2}=20\,\text{м}\ \Rightarrow\ t=\sqrt{\frac{2h}{g}}=2\,\text{с}.$

Мощность равна отношению работы, выполняемой за некоторый промежуток времени:


$P=\frac At=\frac{mgh}{t}=100\,\text{Вт}.$

Ответ: $100$.
}

% Source sdamgia-27081: https://phys-ege.sdamgia.ru/problem?id=27081
\BankTask{sdamgia-27081}{
Покоившееся тело массой $4\,\text{кг}$ начало двигаться вдоль прямой под действием силы, которая в течение $5$ секунд развивала среднюю мощность $10\,\text{Вт}$. Какую скорость в результате приобрело это тело, если другие силы работы не совершали?
}{
\smallskip\noindent\textbf{Решение}\par
Средняя мощность, развиваемая силой, равна $N_{\text{ср}}=Fv_{\text{ср}}$ причем при равноускоренном движении средняя скорость равняется $v_{\text{ср}}=\frac{v+v_0}2=\frac v2$ По второму закону Ньютона $F=ma$, где ускорение $a=\frac{v-v_0}{t}=\frac vt$ Объединяя выражения, получаем значение приобретенной скорости:


$v=\sqrt{\frac{2N_{\text{ср}}t}{m}}=\sqrt{\frac{2\cdot10\,\text{Вт}\cdot5\,\text{с}}{4\,\text{кг}}}=5\,\frac{\text{м}}{\text{с}}.$

Ответ: $5$.
}

% Source sdamgia-23290: https://phys-ege.sdamgia.ru/problem?id=23290
\BankTask{sdamgia-23290}{
Подъемный кран равномерно поднимает груз массой $1000\,\text{кг}$. В таблице приведена зависимость высоты $h$ этого груза над землей от времени подъема $t$. Какую мощность развивает кран при поднятии груза?
 
\begin{center}\begin{tabular}{|c|c|c|c|c|}\hline
$h$, м&$2$&$4$&$6$&$8$\\\hline
$t$, с&$5$&$10$&$15$&$20$\\\hline
\end{tabular}\end{center}
}{
\smallskip\noindent\textbf{Решение}\par
Исходя из табличных данных, можно сделать вывод, что груз поднимается равномерно. Значит, скорость тела равна $v=\frac{\Delta h}{\Delta t}=\frac{4\,\text{м}-2\,\text{м}}{10\,\text{с}-5\,\text{с}}=0{,}4\,\text{м/с}$ При равномерном движении сила тяги двигателя равна силе тяжести, а мощность рассчитывается по формуле: $N=F_{\text{Т}}v=mgv=1000\,\text{кг}\cdot10\,\text{м/с}^2\cdot0{,}4\,\text{м/с}=4000\,\text{Вт}=4\,\text{кВт}.$
 
Ответ: $4$.
}

% Source sdamgia-32749: https://phys-ege.sdamgia.ru/problem?id=32749
\BankTask{sdamgia-32749}{
Вьючный осел, идущий по горной тропе, поднимает тюк массой $30\,\text{кг}$ на высоту $120\,\text{м}$ за $2$ часа. Какую среднюю полезную мощность развивает при этом животное?
}{
\smallskip\noindent\textbf{Решение}\par
При подъеме тюка сила, которую прикладывает осел, равна силе тяжести тюка. Средняя скорость при подъеме равна: 
$v_{\text{ср}}=\frac ht$ Тогда средняя мощность, развиваемая ослом, равна:


$N_{\text{ср}}=F\cdot v_{\text{ср}}=\frac{mgh}t=\frac{30\cdot10\cdot120}{7200}=5\,\text{Вт}.$

Ответ: $5$.
}

% Source sdamgia-7280: https://phys-ege.sdamgia.ru/problem?id=7280
\BankTask{sdamgia-7280}{
Брусок массой $5\,\text{кг}$ равномерно перемещают по горизонтальной поверхности со скоростью $1\,\text{м/с}$, прикладывая к нему постоянную силу $4\,\text{Н}$, направленную под углом $60^\circ$ к горизонту. Чему равна мощность силы $F$?
\begin{center}\begin{tikzpicture}[task graph]
\draw[task curve] (-1.2,0)--(3.2,0);
\foreach \x in {-1.2,-.95,...,3.1}{\draw[task tick] (\x,0)--++(-.18,-.18);}
\draw[task curve,fill=orange!30] (0,0) rectangle (1.5,1);
\fill (1.5,.5) circle (1.5pt);
\draw[task axis] (1.5,.5)--(3,.5);
\draw[task curve,->] (1.5,.5)--++(60:2) node[above left] {$\vec F$};
\draw[task curve] (2.05,.5) arc (0:60:.55);
\node at (2.55,1.1) {$60^\circ$};
\end{tikzpicture}\end{center}
}{
\smallskip\noindent\textbf{Решение}\par
Формула для вычисления мощности силы: $N=FV\cos\alpha=4\cdot1\cdot0{,}5=2\,\text{Вт}.$
 
Ответ: $2$.
}

% Source sdamgia-29128: https://phys-ege.sdamgia.ru/problem?id=29128
\BankTask{sdamgia-29128}{
В велотренажерах для регулировки физической нагрузки тренирующихся на них спортсменов в настоящее время часто используются электродинамические тормозящие устройства, позволяющие плавно регулировать усилия, необходимые для вращения педалей с определенной скоростью. Вращение от педалей передается на массивный токопроводящий диск, находящийся между двумя сильными неподвижными магнитами, расстояние от которых до диска можно регулировать. Взаимодействие возникающих в диске индукционных токов с магнитами тормозит вращение диска, а, следовательно, и педалей, заставляя прикладывать к ним регулируемые по величине силы.
Пусть спортсмен крутит педали, находящиеся на расстоянии $R$ = $20\,\text{см}$ от их оси вращения, с частотой $\nu$ = $15$ оборотов в минуту, прикладывая к каждой из педалей в направлении ее движения постоянную по модулю вращающую силу $F$ = $50\,\text{Н}$. На сколько градусов нагреется алюминиевый диск массой $m$ = $5\,\text{кг}$ за время $t$ = $30$ минут работы в таком режиме? Считайте, что вся работа спортсмена расходуется только на равномерный разогрев диска.
}{
\smallskip\noindent\textbf{Решение}\par
$1$. Найдем вначале мощность $P$, развиваемую спортсменом во время тренировки на велотренажере. Поскольку сила $F$ прикладывается в направлении движения каждой из двух педалей по окружности радиусом $R$ с линейной скоростью $v=\omega R$ где угловая скорость $\omega$ связана с частотой $\nu$ вращения педалей соотношением $\omega=2\pi\nu$ то: 


$P=F\cdot2v=F\cdot2\omega R=4\pi\nu FR$

$2$. Работа, совершенная спортсменом за время $t$, равна: 


$A=P\cdot t=4\pi\nu FRt$

$3$. Поскольку по условию вся эта работа превращается в количество теплоты $Q$, идущее на нагревание алюминиевого диска массой $m$ с удельной теплоемкостью $c$ на $\Delta T$ градусов, то: 


$A=Q=cm\Delta T=4\pi\nu FRt$

$4$. Таким образом, из последнего уравнения получаем: $\Delta T=\frac{4\pi\nu FRt}{cm}$
$5$. Переводя численные данные из условия в систему СИ, используя табличные данные и подставляя их все в полученное выражение, получаем: 


$\Delta T=\frac{4\pi\nu FRt}{cm}=\frac{4\cdot3{,}14\cdot0{,}25\cdot50\cdot0{,}2\cdot1800}{900\cdot5}=12{,}56\,{}^\circ\text{C}.$

Ответ: $12{,}56\,{}^\circ\text{C}.$
}

% Source sdamgia-7718: https://phys-ege.sdamgia.ru/problem?id=7718
\BankTask{sdamgia-7718}{
Известно, что «лошадиная сила» (л. с.) равна мощности $75\,\text{кгс}\cdot\text{м/с}\approx735\,\text{Вт}$, а средний человек при длительной работе развивает мощность около $0{,}16\,\text{л.\,с.}$ и кратковременно может превышать это ограничение. Человек, стараясь после отключения электричества в сети осветить свое жилище, используя электрогенератор с механическим приводом с КПД $\eta$ = $60\%$, вращает ротор генератора через редуктор за ручку, находящуюся на расстоянии $R$ = $0{,}5\,\text{м}$ от оси, со скоростью $n=20\,\text{об/мин}$, прикладывая к ручке силу $F$ = $100\,\text{Н}$. Сможет ли он долго поддерживать горение лампочки мощностью $P$ = $60\,\text{Вт}$, и не перегорит ли она от перенапряжения (лампочка рассчитана на номинальное напряжение $220\,\text{В}$, но не более $235\,\text{В}$, а напряжение генератора прямо пропорционально скорости вращения ротора)?
}{
\smallskip\noindent\textbf{Решение}\par
КПД генератора с механическим приводом равен отношению его электрической мощности к механической мощности, развиваемой человеком:


$\eta=\frac{P_{\text{эл}}}{P_{\text{мех}}}$

Механическая мощность, развиваемая человеком, равна $P_{\text{мех}}=Fv$ где скорость движения по окружности $v=2\pi Rn$ Откуда получаем:


$P_{\text{мех}}=F\cdot2\pi n\cdot R\approx100\cdot6{,}28\cdot\frac{20}{60}\cdot0{,}5\approx104{,}7\,\text{Вт}$

а электрическая мощность генератора $P_{\text{эл}}=\eta P_{\text{мех}}\approx62{,}8\,\text{Вт}$ что достаточно для питания лампочки мощностью $60\,\text{Вт}$. 
Напряжение питания по условию прямо пропорционально скорости вращения


$220\cdot\sqrt{\frac{62{,}8}{60}}\approx225\,\text{В}$

так что лампочка еще не перегорит.
Средняя механическая мощность человека при длительной работе равна по условию $0{,}16\,\text{л.\,с.}\approx117{,}6\,\text{Вт}$, так что человек сможет долго освещать свое жилище, не перенапрягаясь.
 
Ответ: Сможет, и лампочка не перегорит.
}

\end{document}
